% Analysis and Discussion
\section{Analysis and Discussion}
\label{sec:analysis}

\subsection{Out-of-Domain Improvement: Observations and Hypotheses}

\subsubsection{Observed Facts}

\begin{itemize}[nosep]
    \item Consistent improvements across HumanEval (+2.86\%), MBPP (+3.00\%), and GSM8K (+1.07\%)
    \item Transfer effect observed using only 1.4\% of available training data (50K/3.6M samples)
    \item Improvements are more pronounced at higher Pass@k values (Pass@10, Pass@20)
\end{itemize}

\subsubsection{Hypotheses (Unverified)}

We propose several hypotheses that may explain the out-of-domain improvements:

\begin{enumerate}[nosep]
    \item \textbf{Short sequence suitability}: Out-of-domain benchmarks feature shorter, function-level code where value changes may better capture token importance. The critic's 63\% accuracy may be sufficient for these simpler structures.

    \item \textbf{Critical token focus}: Weighted cross-entropy may allocate larger gradients to ``decisive tokens'' (function signatures, return statements, boundary conditions) that disproportionately affect correctness.

    \item \textbf{Regularization effect}: Token weighting may act as implicit regularization, preventing overfitting to less important tokens in the training data.
\end{enumerate}

\subsubsection{Required Verification}

To validate these hypotheses, we would need:
\begin{itemize}[nosep]
    \item Token-level weight distribution analysis by code token type (keywords, identifiers, operators)
    \item Attention and gradient analysis to confirm ``important token'' hypothesis
    \item \textbf{Random weight baseline comparison} (critical missing experiment)
\end{itemize}

\subsection{In-Domain (CodeContests) Decline: Cause Analysis}

\subsubsection{Observed Facts}

\begin{itemize}[nosep]
    \item Average decline of -1.05\% across Pass@k values
    \item Most pronounced at Pass@5 (-1.75\%) and Pass@10 (-1.97\%)
    \item Pass@20 shows no difference (both at 10.00\%)
\end{itemize}

\subsubsection{Hypotheses (Prioritized)}

We propose the following causes, ordered by estimated likelihood:

\begin{enumerate}
    \item \textbf{Critic quality limitation (High confidence)}:
    The critic achieves only 63.1\% validation accuracy. For long competitive programming solutions (often $>$1000 tokens), accumulated errors in token-level value estimates may produce misleading advantage signals.

    \item \textbf{GAE credit assignment breakdown (Medium confidence)}:
    GAE with $\lambda=0.95$ accumulates TD-errors over the sequence. For very long sequences, this accumulation may become unstable, assigning incorrect credit to early tokens based on later value fluctuations.

    \item \textbf{Data coverage insufficiency (Medium confidence)}:
    Using only 1.4\% of available data may fail to cover the difficulty distribution of competitive programming problems adequately. The training distribution may not represent the evaluation distribution.

    \item \textbf{Critic backbone capacity (Lower confidence)}:
    The 2.7B critic model may lack capacity to accurately estimate values for solutions generated by a 7B policy model, creating a capability mismatch.
\end{enumerate}

\subsection{Critic Quality and Downstream Performance (Unverified)}
\label{sec:critic_quality}

\subsubsection{Current State (Facts Only)}

\begin{itemize}[nosep]
    \item Critic validation accuracy: 63.1\%
    \item Despite this moderate accuracy, out-of-domain improvements are observed
    \item Pairwise ranking outperformed MSE regression (55\% $\rightarrow$ 63\%)
\end{itemize}

\subsubsection{Unresolved Questions}

\begin{enumerate}
    \item \textbf{Is 63\% accuracy sufficient?} Can a critic with only 63\% pairwise accuracy provide ``meaningful'' token-level weights, or are the observed improvements due to other factors?

    \item \textbf{Information vs. weighting effect}: How much improvement comes from the \textit{information} in the weights (correct importance signals) vs. the \textit{weighting mechanism itself} (any non-uniform weighting might help)?

    \item \textbf{Accuracy-performance relationship}: Would increasing critic accuracy to 70\%, 80\%, or higher yield proportionally better downstream performance?
\end{enumerate}

\subsubsection{Required Future Experiments}

\begin{itemize}[nosep]
    \item \textbf{Random/shuffled weight A/B test}: Train policy with random weights sampled from the same distribution as value-based weights. If performance is similar, the improvement may be due to weighting mechanism rather than value information.

    \item \textbf{Critic checkpoint sweep}: Use critic checkpoints at different training stages (varying accuracy levels) to train identical policies, measuring the accuracy-performance correlation.
\end{itemize}

\subsection{Data Efficiency Observations}

A notable finding is that WMTP achieves improvements using only \textbf{1.4\% of available training data}:

\begin{itemize}[nosep]
    \item Total available: $\sim$3.6M samples
    \item Actually used: 50K samples
    \item Usage ratio: 1.4\%
\end{itemize}

This suggests that token weighting may be particularly effective in \textbf{low-data regimes}, where focusing learning on important tokens could compensate for limited exposure to diverse examples. However, this observation requires controlled experiments varying data quantity to confirm.

\subsection{Summary of Analysis}

\begin{table}[h]
\centering
\caption{Summary of analysis: observations, hypotheses, and verification status.}
\label{tab:analysis_summary}
\begin{tabular}{lp{6cm}c}
\toprule
\textbf{Finding} & \textbf{Explanation} & \textbf{Verified?} \\
\midrule
Out-of-domain $\uparrow$ & Short sequences + token focus & No (hypothesis) \\
In-domain $\downarrow$ & Critic quality + GAE breakdown & No (hypothesis) \\
63\% critic $\rightarrow$ improvement & Sufficient for simple tasks? & No (needs baseline) \\
1.4\% data $\rightarrow$ transfer & Data-efficient weighting? & No (needs sweep) \\
\bottomrule
\end{tabular}
\end{table}
